% =====================================================================
%  一、问题背景与重述
% =====================================================================
\section{问题背景与重述}

\subsection{问题背景}
随着无线电技术的普及和广泛应用，频谱管理中的干扰源定位与清除问题日益突出。
干扰不仅影响正常通信，还可能危及公共安全。传统清除方式通常受环境遮挡、低功率发射等影响，
精准定位耗时长，且危险环境下人工作业风险高，因此研发机器狗自动定位清除系统代替人工执行任务具有现实需求。
但任务中通常面临着目标区域范围较大，干扰源数量、频道、位置、有效接收半径及类型未知，
测向存在误差，移动、检测、切换频道和清除过程存在耗时，定向干扰源的发射方向无法辨识等困难。
如何在信息不完全且时间受限的条件下高效完成任务，是该自动定位清除系统亟待解决的问题。

\subsection{问题重述}
考虑到上述实际情况中面临的困难，本文需依次解决以下四个问题：

\begin{enumerate}
    \item 已知若干检测点的位置及干扰源关于这些检测点的示向度，给出用交会定位法计算
          可能位置区域（多边形）直径的算法，并回答以该直径为直径的圆能否覆盖此区域。
    \item 已知在某一个检测点处测得的示向度，给出第二个检测点的选择策略，并指出该点
          可取在什么范围（候选区域）之内。
    \item 区域内存在多个干扰源，它们占用互不相同的频道且总个数未知，机器狗从给定
          起点出发自动定位并逐个清除。要求给出策略与算法，经模拟器演练检验后给出
          清除比例与平均定位清除时间，并据以完成正式测试。
    \item 区域内除全向干扰源外还混有定向干扰源，其个数与发射方向均未知，其余条件与
          问题三相同。要求给出相应的检测与清除策略及算法。
\end{enumerate}
